% Uvod
%
%
\phantomsection
\addcontentsline{toc}{chapter}{Uvod}
\chapter*{Uvod}
\label{chap:intro}
%
%
Serijska komunikacija smatra se jednim od najstarijih i
najrasprostranjenijih načina razmjene podataka između mikrokontrolera i
perifernih uređaja, a \foreignlanguage{english}{UART}
(\foreignlanguage{english}{Universal Asynchronous Receiver-Transmitter})
protokol je zbog svoje jednostavnosti i malog broja potrebnih linija
ostao u širokoj upotrebi i danas, uprkos postojanju brojnih novijih
protokola. Posebno zanimljivim svojstvom ovog protokola smatra se
njegova asinhronost, odnosno odsustvo zajedničke linije takta između
predajnika i prijemnika, zbog čega se tačan trenutak očitavanja svakog
pojedinačnog bita mora unaprijed izračunati na osnovu poznate brzine
prijenosa. Format okvira, koji se sastoji od tačno određenog niza stanja
kroz koja linija prolazi u tačno određenom vremenskom rasporedu, smatra
se pogodnim praktičnim primjerom konačnog automata čije prelaze između
stanja ne pokreće samo promjena ulaznog signala, već i protok vremena.

Iako su savremeni mikrokontroleri, uključujući i
\foreignlanguage{english}{ATmega328P} korišten u ovom radu, opremljeni
ugrađenim hardverskim \foreignlanguage{english}{UART} perifernim
modulom kojim se ova komunikacija izvodi automatski, kod softverske
realizacije istog protokola, poznate kao \foreignlanguage{english}{bit-banging},
tajming svake pojedinačne operacije mora se precizno kontrolisati na
nivou broja ciklusa procesora. Zbog ove osobine, softverska realizacija
smatra se pogodnom za razumijevanje veze između arhitekture procesora,
skupa njegovih instrukcija i vremenski osjetljivih protokola opisanih
formalnim automatima. Svrha ovog rada je da se na konkretnom primjeru
pokaže postupak kojim se, polazeći od poznatog broja ciklusa svake
korištene instrukcije, konstruiše program koji ostvaruje zadatu brzinu
prijenosa isključivo korištenjem digitalnih ulazno-izlaznih pinova.

U radu će se odgovoriti na pitanje kako se, bez upotrebe hardverskog
\foreignlanguage{english}{UART} modula, može ostvariti tačno određena
brzina prijenosa korištenjem petlje za kašnjenje čiji je broj ponavljanja
izveden iz broja ciklusa procesora. Razmotrit će se koje su instrukcije
\foreignlanguage{english}{AVR} arhitekture najpogodnije za rad sa
pojedinačnim bitom ulazno-izlaznog registra i na koji način njihov izbor
utiče na preciznost i predvidljivost tajminga. Poseban naglasak bit će
stavljen na pitanje kako se ista petlja za kašnjenje, uz odgovarajuće
prilagođavanje okolnog koda, može iskoristiti i za slanje i za prijem
podataka, uprkos tome što se struktura koda oko poziva te petlje
razlikuje u oba slučaja.

Rad se sastoji iz tri poglavlja. U prvom poglavlju izlaže se teorijska
osnova neophodna za razumijevanje ostatka rada. Opisuje se format
\foreignlanguage{english}{UART} okvira i način izvođenja brzine
prijenosa, kao i razlog očitavanja na sredini bita, a zatim se uvode
dijelovi \foreignlanguage{english}{AVR} arhitekture koji se koriste u
implementaciji, uključujući registre opće namjene, status registar
\foreignlanguage{english}{SREG}, ulazno-izlazne portove i trajanje
izvršavanja instrukcija u ciklusima.

U drugom poglavlju izvodi se vremenska analiza programa na nivou
pojedinačnih instrukcija. Opisuje se skup korištenih instrukcija zajedno
sa njihovim opkodom i tačnim brojem ciklusa izvršavanja, a zatim se, na
osnovu tako utvrđenih vrijednosti, izvodi opšta formula kojom se
određuje broj ponavljanja petlje za kašnjenje. Formula se potom
primjenjuje na dva konkretna slučaja, punu bit periodu i poravnanje na
sredinu bita prilikom prijema.

U trećem poglavlju opisuje se praktična implementacija programa.
Objašnjava se sintaksa i skup direktiva korištenih pri pisanju izvornog
koda, zatim algoritam slanja i prijema podataka uz obrazloženje odabira
korištenih instrukcija i pomoćnih makroa za usklađivanje broja ciklusa.
Na kraju poglavlja opisuje se razvojno okruženje i postupak testiranja
realizovanog programa na fizičkoj razvojnoj ploči.

Na kraju rada nalazi se zaključak, u kojem su sistematizovani
najvažniji rezultati i doprinosi rada kao cjeline. Pored zaključka, u
radu se nalaze i popis slika, kao i lista korištenih referenci.

\newpage
